\par\SourceLeaf{140}
Da dieses Gutachten und die ihm gegebene Auslegung die eigentliche Grundlage von derjenigen Einrichtung bilden, die den Hauptgegenstand der 1. Abtheilung meiner Denkschrift ausmacht, nemlich von der Progressionsnormirung, so werde ich solches, weiter unten, einer ausführlichen und genauen Prüfung unterwerfen, und beschränke mich daher hier, einstweilen, auf folgende Bemerkungen.

Aus den Acten ist nicht zu ersehen, weshalb \textsc{Kritter} im Anfange seines Gutachtens von 22—24 verehelichten Professoren spricht, da P. in seinem Schreiben ausdrücklich 26, und nur diese Zahl, genannt hatte. Ich vermuthe aber, dass \textsc{Kritter} jene Zahlen 22—24 wie die für eine \textit{frühere} Zeit gültigen angenommen hat. Meines Wissens sind aber keine vollständige Register über die personlichen Verhältnisse der Witwenkassen-Mitglieder in der Art geführt, dass für jeden beliebigen Zeitpunkt die Anzahl der \textit{Verehelichten} daraus entnommen werden könnte. Entweder also hat \textsc{Kritter} jene Zahlen nur aus der Zahl aller Participanten in früherer Zeit nach einer ungefähren Schätzung \textit{geschlossen}, oder sie beruhen auf besondern Mittheilungen, welche dann, der Natur der Sache nach, sich nur auf Zeitpunkte beziehen können, die nicht viele Jahre rückwärts lagen, und im Gedächtnisse noch fortlebten.

\textsc{Kritter}’s Antwort auf die erste Frage besteht dann kurz darin, dass man für die \textit{Zeit}, wo das Maximum eingetreten sei, Eine Witwe gegen etwa zwei stehende Ehen rechnen könne, also für jene 22—24 Ehen 12 Witwen, was sich wie er angibt nach dem Durchschnitt der letzten 17 Jahre in so fern bestätigt habe, als bald mehr bald weniger als 12 Witwen vorhanden gewesen seien. Dass dann die der dermaligen Zahl von 26 Ehen entsprechende Witwenzahl um Eine grösser sein würde, ist nicht ausdrücklich gesagt, aber implicite darin enthalten. Auf die zweite Frage gibt er an, dass nach den Erfahrungen der Bremischen Witwenkasse, wo die Waisenpensionirung nur bis zum \(18^{\text{ten}}\) Jahre daure, man auf eine Vergrösserung der Pensionenzahl um den sechsten Theil rechne; bei der hiesigen also, wo die Dauer 2 Jahre länger sein solle, doch wohl etwas mehr annehmen müsse. — Bei den 24 Ehen kommen wir demnach auf etwas mehr als 14, bei den 26 Ehen, nachdem ihre Wirkung ganz eingetreten, auf 15 nach \textsc{Kritter}’s Worten, oder auf etwas mehr als \(15\tfrac12\), d. i. auf nahe 16 nach den von ihm ausgesprochenen Grundsätzen.

Ob, \textit{ganz abgesehen von der nähern Prüfung des Inhalts des Gutachtens}, eine derartige Behandlung des Gegenstandes eine \textit{umständliche} und \textit{genaue} Untersuchung, wie \textsc{Kaestner} für nothwendig gehalten hatte, genannt werden könne, lasse ich hier auf sich beruhen. P. entwarf nun aber, auf den Grund dieses Gutachtens, ein P. M., worin er zeigt, dass wenn die von dem Curatorium vorgeschlagene Erhöhung der Pensionen, ohne weitere besondere Vergrösserung für die sechs ältesten Witwen, für alle gleichmässig auf 130 Rth. festgesetzt werde, dies ohne alle Gefahr für die Kasse auch dann geschehen könne, wenn die jährlichen Beiträge nicht erhöhet würden; dass aber, im Fall die Erhöhung der Beiträge auf die vom Curatorium angegebene Art, angenommen werde, auch die Verlängerung der Waisenpensionen um so sicherer eingeführt werden könne, weil nach den obwaltenden Umständen ein baldiges Wirksamwerden dieser Abänderung nicht zu erwarten sei. P. schliesst endlich seinen Vortrag, dem ich, \textit{bis hieher}, meinen vollen Beifall zu geben keinen Anstand nehme, mit folgendem kurzen Zusatz, den ich, da hier \textit{zum erstenmale} der Gegenstand meiner eignen Untersuchung, nemlich die Progressionsnormirung, auf den Schauplatz tritt, vollständig und treu mit P.’s eignen Worten hier abschreibe:

\begin{quote}
‘Bei der allgemeinen Erhöhung der Pensionen auf 130 Rth. scheint mir nicht die mindeste Gefahr zu sein:
\end{quote}
